\documentclass[aspectratio=169,10pt,xcolor={dvipsnames}]{beamer}

%\includeonlyframes{current}

\usetheme{blei} % at ~/texmf/tex/latex/local/
\addbibresource{library.bib}
\usetikzlibrary{calc}

% https://tex.stackexchange.com/a/1475/193011
% PGF doesn't use TikZ, but the tikzplotlib backend does
%\usetikzlibrary{external}
%\tikzexternalize[prefix=tikz/]

%%%%%%%%%%
% \notes %
%%%%%%%%%%
\usepackage{pgfpages}
%\setbeameroption{hide notes} % Only slides
%\setbeameroption{show only notes} % Only notes
%\setbeameroption{show notes on second screen=right} % Both
\setbeamertemplate{note page}{\pagecolor{yellow!5}\vfill\insertnote\vfill}

% https://latex.org/forum/viewtopic.php?t=11357
\usepackage{xparse}
\NewDocumentCommand\notes{>{\SplitList{;}}m}
{
	\note{%
    \begin{itemize}
      \ProcessList{#1}{ \insertitem }
    \end{itemize}
    }
}
\newcommand\insertitem[1]{\item #1}


%%%%%%%%%%%%
% Packages %
%%%%%%%%%%%%
\usepackage{amsmath}
\usepackage{physics}
\usepackage{pgf}
\DeclareUnicodeCharacter{2212}{-} % Needed for matplotlib's pgf output
\usepackage{pythonhighlight}
\usepackage{subfig}
\usepackage{marvosym}
\usepackage[normalem]{ulem}
\let\Cross\relax
\usepackage{bbding}
\usepackage{ragged2e}
\usepackage{multicol}

%\usepackage{bm} % Clashes with main font, so use \boldsymbol

%%%%%%%%%%%%
% Commands %
%%%%%%%%%%%%
\newcommand{\xmax}{x_{\text{max}}}
\newcommand{\Kmax}{K_{\text{max}}}

\newcommand{\bx}{{\bm x}}
\newcommand{\bxnull}{{\bm x_0}}
\newcommand{\blambda}{{\bm \lambda}}
\newcommand{\ba}{{\bm a}}
\newcommand{\bb}{{\bm b}}
\newcommand{\btheta}{{\bm \theta}}
\newcommand{\balpha}{{\bm \alpha}}
\newcommand{\bu}{{\bm u}}
\newcommand{\bd}{{\bm d}}
\newcommand{\bA}{{\bm A}}
\newcommand{\bbeta}{{\bm \beta}}

\newcommand{\barx}[1]{\overline{x_{#1}}}
\newcommand{\baru}[1]{\overline{u_{#1}}}
\newcommand{\bbarx}{\overline{\bm x}}
\newcommand{\bbarxnull}{\overline{\bm{x_0}}}
\newcommand{\bbaru}{\overline{{\bm u}}}

\DeclareMathOperator{\Pareto}{\mathsf{Pareto}}
\DeclareMathOperator{\ParetoChain}{\mathsf{ParetoChain}}
\DeclareMathOperator{\Jeffreys}{\mathsf{Jeffreys}}
\DeclareMathOperator{\Exp}{\mathsf{Exp}}

\newcommand{\priorinformation}{\mathcal{I}}
\newcommand{\likelihood}{\mathcal{L}}
\newcommand{\model}{\mathcal{M}}

\newcommand{\orderedregion}{\mathcal{R}}

\newcommand{\data}[1]{{\mathsf{#1}}}
\newcommand{\little}{\data{little}}
\newcommand{\until}{\data{until}}
\newcommand{\shore}{\data{shore}}
\newcommand{\that}{\data{that}}
\newcommand{\you}{\data{you}}

\newcommand{\numb}[1]{\mathsf{#1}}
\newcommand{\bm}{\boldsymbol}

\let\originalincludegraphics\includegraphics
\newcommand{\missingfigure}[1]{%
\fbox{\parbox[c][2cm][c]{4cm}{\centering\scriptsize Missing figure\\\texttt{#1}}}%
}
\renewcommand{\includegraphics}[2][]{%
\IfFileExists{#2}{\originalincludegraphics[#1]{#2}}{\missingfigure{#2}}%
}
\newcommand{\pgffigure}[2]{%
\IfFileExists{img/#1}{\scalebox{#2}{\input{img/#1}}}{\scalebox{#2}{\missingfigure{img/#1}}}%
}

\newcommand{\citelink}[2]{\href{#1}{\small\color{gray}[#2]}}

\definecolor{mplblue}{HTML}{1f77b4}
\definecolor{mplorange}{HTML}{ff7f0e}

% https://tex.stackexchange.com/a/326380/193011
\usepackage{xcolor}
% Syntax: \colorboxed[<color model>]{<color specification>}{<math formula>}
\newcommand*{\colorboxed}{}
\def\colorboxed#1#{%
  \colorboxedAux{#1}%
}
\newcommand*{\colorboxedAux}[3]{%
  % #1: optional argument for color model
  % #2: color specification
  % #3: formula
  \begingroup
    \colorlet{cb@saved}{.}%
    \color#1{#2}%
    \boxed{%
      \color{cb@saved}%
      #3%
    }%
  \endgroup
}

%%%%%%%%%%%%%%%%%%%%%%%%%%%%%%%%%%%%%%%%%%%%%%%%%%%%%%%%%%%%%%%%%%%%%%%%%%%%%%%%
% Content-slide layout (logos, section label, page number + corner arc)      %
%   - bottom left  : uni.lu short logo                                        %
%   - bottom right : SnT logo                                                 %
%   - top left     : current \section name (small, grey)                      %
%   - top right    : frame number inside a light corner arc                   %
%   - frame title  : large bold navy                                          %
% Plain frames (e.g. the title slide) suppress the headline/footline, so      %
% these decorations appear only on ordinary content frames.                   %
%%%%%%%%%%%%%%%%%%%%%%%%%%%%%%%%%%%%%%%%%%%%%%%%%%%%%%%%%%%%%%%%%%%%%%%%%%%%%%%%

\definecolor{titlenavy}{HTML}{1d2c44}   % big frame title
\definecolor{sectiongrey}{HTML}{8a93a3} % top-left section label
\definecolor{arcgrey}{HTML}{c9cedb}     % top-right corner arc

\setbeamercolor{background canvas}{bg=white}

% ── Frame title: large, bold, navy ───────────────────────────────────────
\setbeamercolor{frametitle}{fg=titlenavy}
\setbeamerfont{frametitle}{series=\bfseries,size=\huge}
\setbeamertemplate{frametitle}{%
  \vspace{0.9em}%
  {\usebeamercolor[fg]{frametitle}\usebeamerfont{frametitle}\insertframetitle\par}%
  \vspace{0.15em}%
}

% ── Headline: section label (top-left) + frame number in a corner arc ─────
%  No on-slide navigation bar — jumping around is done via the bottom-centre
%  outline button and the Contents/divider pages.
\setbeamertemplate{headline}{%
  \begin{tikzpicture}[remember picture, overlay]
    % small light quarter-circle, detached from both edges
    \draw[arcgrey, line width=0.9pt, line cap=round]
      ([xshift=-1.05cm, yshift=-0.40cm]current page.north east)
      arc[start angle=180, end angle=270, radius=0.75cm];
    % frame counter, at the arc's centre (current/total)
    \node[anchor=center, text=titlenavy, font=\bfseries\scriptsize]
      at ([xshift=-0.40cm, yshift=-0.50cm]current page.north east)
      {\insertframenumber/\inserttotalframenumber};
    % section label, top-left (context only — not a nav control)
    \node[anchor=north west, text=sectiongrey, font=\small]
      at ([xshift=0.045\paperwidth, yshift=-0.30cm]current page.north west)
      {\insertsection};
  \end{tikzpicture}%
}

% ── Table-of-contents styling (master outline + section dividers) ─────────
\setbeamercolor{section in toc}{fg=titlenavy}
\setbeamercolor{section in toc shaded}{fg=titlenavy!68!white}
\setbeamercolor{subsection in toc}{fg=titlenavy}
\setbeamercolor{subsection in toc shaded}{fg=titlenavy!58!white}
\setbeamerfont{section in toc}{series=\bfseries}
\setbeamerfont{subsection in toc}{size=\footnotesize}

% ── Outline button: small clickable "contents" icon → current section's
%    divider (each divider lists every section, so it is a full nav hub) ────
\newcommand{\outlinebutton}{%
  \hyperlink{sec\arabic{section}}{%
    \scalebox{0.7}{%
      
\begin{tikzpicture}[x=1mm, y=1mm, baseline=-0.4mm]
        \foreach \yy in {0,-1.6,-3.2}{%
          \fill[sectiongrey] (0,\yy) circle (0.42);
          \draw[sectiongrey, line width=0.8pt, line cap=round] (1.4,\yy) -- (5.4,\yy);
        }%
      \end{tikzpicture}%
    }%
  }%
}

% ── Footline: uni.lu (left) · outline button (centre) · SnT (right) ───────
\setbeamertemplate{footline}{%
  \hbox to\paperwidth{%
    \hspace*{0.03\paperwidth}%
    \raisebox{0pt}{
\includegraphics[height=0.055\paperheight]{Assets/UniLu_short.png}}%
    \hfil
    \raisebox{0.013\paperheight}{\outlinebutton}%
    \hfil
    \raisebox{0pt}{
\includegraphics[height=0.050\paperheight]{Assets/SnT_logo.png}}%
    \hspace*{0.03\paperwidth}%
  }%
  \vspace*{0.025\paperheight}%
}

% ── Blocks: navy title bar, light body, rounded ───────────────────────────
\setbeamertemplate{blocks}[rounded][shadow=false]
\setbeamercolor{block title}{bg=titlenavy,fg=white}
\setbeamercolor{block body}{bg=titlenavy!7,fg=black}

% ── Section divider: a clean "Contents" page before each section ──────────
%  Current section highlighted with its subsections; other sections shaded
%  with their subsections hidden. TOC entries are clickable.
\AtBeginSection[]{%
  {%
  \usebackgroundtemplate{%
    \includegraphics[width=\paperwidth,height=\paperheight]{Assets/outline_bg.png}%
  }%
  \begin{frame}[plain,noframenumbering,t]{Contents}
    \hypertarget{sec\arabic{section}}{}%
    \begin{tikzpicture}[remember picture, overlay]
      \node[anchor=south west]
        at ([xshift=0.03\paperwidth, yshift=0.025\paperheight] current page.south west)
        {
\includegraphics[height=0.055\paperheight]{Assets/UniLu_short.png}};
      \node[anchor=south east]
        at ([xshift=-0.03\paperwidth, yshift=0.025\paperheight] current page.south east)
        {
\includegraphics[height=0.050\paperheight]{Assets/SnT_logo.png}};
    \end{tikzpicture}
    \vspace{0.25em}
    {\footnotesize
    \setlength{\columnsep}{1.2cm}
    \begin{multicols}{2}
      \tableofcontents[currentsection,
                       sectionstyle=show/shaded,
                       subsectionstyle=show/show/hide]
    \end{multicols}
    }
  \end{frame}
  }%
}

%%%%%%%%%
% Title %
